% Figure 11 draft A: the stack as on the cover, twelve sheets in six version pairs over V; the two lifts t then b and
% b then t from the sheet X end on different sheets (gold rings); the eight sheets not visited drawn lighter.
\documentclass[10pt,border=1mm]{standalone}
\usepackage[T1]{fontenc}
\usepackage{figurestyle}
\input{primitives.tex}
\input{molecules.tex}
\input{numbers.tex}
% sheet parallelogram: \sheet{x}{y}{style}{fill}  (3.6 wide, as on the cover)
\newcommand\sheet[4]{\path[fill=#4] ({#1-1.8},#2)--({#1+1.8},#2)--({#1+2.5},{#2+.5})--({#1-1.1},{#2+.5})--cycle; \draw[#3] ({#1-1.8},#2)--({#1+1.8},#2)--({#1+2.5},{#2+.5})--({#1-1.1},{#2+.5})--cycle;}
\newcommand\goldring[1]{\draw[line width=1pt,draw=ColGoldStroke] (#1) circle[radius=.16];}

\begin{document}
\begin{tikzpicture}[plate,light/.style={line width=.4pt,draw=PlateInk!55}]
\path[use as bounding box] (0,0) rectangle (15,12.6);
\begin{scope}[shift={(4.2,.9)}]
  \sheet{0}{0}{heavy}{white}
  \coordinate (bx) at (.35,.25); \fill[PlateInk] (bx) circle[radius=1.2pt];
  \draw[tlift] (bx) .. controls (-1.5,.72) and (-1.5,-.22) .. (bx); \draw[blift] (bx) .. controls (2.2,-.22) and (2.2,.72) .. (bx);
  \node[sub,anchor=west,inner sep=1pt] at (2.6,.25) {$V$}; \node[sub,anchor=north,inner sep=2pt] at (bx) {$[X]$};
  \node[sub] at (-1.2,.5) {$\ell_t$}; \node[sub] at (2.05,.5) {$\ell_b$};
  \foreach \k/\ga/\gb/\pl in {0/X/bX/H,1/tX/tbX/tH,2/t^2X/t^2bX/t^2H,3/uX/ubX/uH,4/tuX/tubX/tuH,5/t^2uX/t^2ubX/t^2uH} {
    \pgfmathsetmacro{\ya}{1.3+1.5*\k}\pgfmathsetmacro{\yb}{\ya+.42}
    \ifnum\k<3 \sheet{0}{\ya}{fine}{white} \sheet{0}{\yb}{fine}{white} \else \sheet{0}{\ya}{light}{white} \sheet{0}{\yb}{light}{white} \fi
    \coordinate (pa\k) at (.35,{\ya+.25}); \coordinate (pb\k) at (.35,{\yb+.25});
    \node[sub,anchor=west,inner sep=1pt] at (1.9,{\ya+.1}) {$\ga$}; \node[sub,anchor=west,inner sep=1pt] at (1.9,{\yb+.42}) {$\gb$};
    \draw[fine] (2.75,\ya)--(2.9,\ya)--(2.9,{\yb+.5})--(2.75,{\yb+.5}); \node[lab,anchor=west] at (3.0,{\ya+.46}) {$\pl$};}
  \draw[tlift] (pa0) .. controls (-1.5,{1.3+.25+.5}) and (-1.5,{1.3+1.5+.25-.5}) .. (pa1);
  \draw[blift] (pa1) .. controls (1.6,{1.3+1.5+.25+.08}) and (1.6,{1.3+1.5+.42+.25-.08}) .. (pb1);
  \draw[blift] (pa0) .. controls (1.6,{1.3+.25+.08}) and (1.6,{1.3+.42+.25-.08}) .. (pb0);
  \draw[tlift] (pb0) .. controls (-1.7,{1.3+.42+.25+1.0}) and (-1.7,{1.3+3.0+.42+.25-1.0}) .. (pb2);
  \foreach \k in {0,1,2} {\hatomat{pa\k}{.07}{ColDot} \hatomat{pb\k}{.07}{ColDot}}
  \foreach \k in {3,4,5} {\draw[light,fill=ColDot!40] (pa\k) circle[radius=.07]; \draw[light,fill=ColDot!40] (pb\k) circle[radius=.07];}
  \goldring{pb1} \goldring{pb2}
  \node[sub,anchor=east] at (-1.95,{1.3+1.5+.42+.25}) {$t$ then $b$}; \node[sub,anchor=east] at (-1.95,{1.3+3.0+.42+.25}) {$b$ then $t$};
\end{scope}
% X and bX, two configurations over one point
\begin{scope}[shift={(11.6,8.6)}]
  \begin{scope}\molsolid\mollabelsinside\pic[scale=.5] at (-1.1,.3) {mol3d-mla-X}; \pic[scale=.5] at (1.2,.3) {mol3d-mla-bX};\end{scope}
  \node[lab] at (-1.1,-1.0) {$X$}; \node[lab] at (1.2,-1.0) {$bX$}; \draw[blift] (-.55,-1.0)--(.65,-1.0);
  \node[sub,anchor=north,text width=5.2cm,align=center] at (0,-1.35) {two configurations over the one point $[X]$; no rotation carries one to the other};
\end{scope}
\end{tikzpicture}
\end{document}
